% ===== Document Preamble =====
\documentclass[11pt]{article}
\usepackage[margin=0.5in]{geometry}
\usepackage{amsmath,amssymb}
\usepackage{xcolor}
\usepackage{enumitem}
\usepackage{needspace}
\usepackage{comment}

% ===== Colors =====
\definecolor{AnsRed}{RGB}{170,0,0}
\definecolor{WrongBlue}{RGB}{0,0,160}

% ===== Choice list styling (tight, consistent) =====
\setlist[enumerate]{%
  leftmargin=1.5em,
  itemsep=0.25em,
  topsep=0.25em,
  label=\textbf{\Alph*.}
}

% ===== Style Wrappers =====
\newcommand{\QuestionStyle}[1]{{\large\bfseries\boldmath #1}}
\newcommand{\AnswerStyle}[1]{{\large\bfseries\boldmath #1}}
\newcommand{\ReasonStyle}[1]{{\normalsize #1}}

% ===== Show / Hide Reasoning Toggle =====
\newif\ifshowreason
\showreasontrue      % Instructor version
% \showreasonfalse   % Student version

% ===== Choice Item Macros (Answer + Reasoning) =====
% Usage:
%   \Correct{Answer}{Reasoning}
%   \Wrong{Answer}{Reasoning}

\newcommand{\Correct}[2]{%
  \item \AnswerStyle{\textcolor{AnsRed}{#1}}%
  \ifshowreason\\[-0.2em]
    \ReasonStyle{\textcolor{AnsRed}{#2}}%
  \fi
}

\newcommand{\Wrong}[2]{%
  \item \AnswerStyle{\textcolor{WrongBlue}{#1}}%
  \ifshowreason\\[-0.2em]
    \ReasonStyle{\textcolor{WrongBlue}{#2}}%
  \fi
}

% ===== Question Block Macro =====
% Usage: \QBlock{Q1}{stem...}{<items>}

\newcommand{\QBlock}[3]{%
  \Needspace{8\baselineskip}%
  \noindent
  \begin{minipage}{\linewidth}
    \QuestionStyle{#1.~#2}

    \vspace{0.25em}
    \begin{enumerate}
      #3
    \end{enumerate}

    \vspace{0.25em}
    \hrule height 0.4pt
  \end{minipage}
  \vspace{0.75em}
}

% ===== Document Start =====

\begin{document}

\begin{center}
{\huge \textbf{Physics Questions}}\\[0.25em]
{\Large Projectile Motion and Spacecraft Dynamics}\\[0.25em]
{\large Multiple choice questions generated below. Answers in RED and incorrect answers in BLUE.}
\end{center}

\vspace{0.75em}

\QBlock{Q1}{A Chinook salmon leaps from the water at a speed of $v_0$ with an angle of $55^\circ$. What is the fish's speed at its maximum height?}{
  \Wrong{$v_0 \sin 55^\circ$}{This choice represents the initial vertical velocity component, which becomes zero at maximum height.}
  \Wrong{0}{The vertical component becomes zero, but the horizontal component remains at maximum height.}
  \Wrong{$v_0$}{The speed at maximum height involves only the horizontal component, which is $v_0 \cos \theta$.}
  \Correct{$v_0 \cos 55^\circ$}{The horizontal velocity, $v_0 \cos \theta$, remains constant throughout the motion since there's no horizontal acceleration.}
}

\QBlock{Q2}{A spacecraft moves with an initial velocity $\vec{v}_0 = (22 \text{ m/s})\hat{i}$, and its boosters provide an acceleration $\vec{a} = (2.5\text{ m/s}^2)\hat{i} + (3.5\text{ m/s}^2)\hat{j}$. Which statement correctly describes the motion?}{
  \Wrong{The y-component of the velocity does not change.}{Since there is acceleration in the $\hat{j}$ direction, the y-component does change.}
  \Wrong{The x-component of the velocity does not change.}{There is an acceleration in the $\hat{i}$ direction, causing the x-component to change.}
  \Wrong{The speed increases but the direction remains unchanged.}{Both speed and direction change because there are accelerations in both x and y directions.}
  \Correct{Both the speed and direction change.}{Accelerations in both x and y cause changes in both speed and direction, altering the spacecraft's trajectory.}
}

\end{document}
